\section*{الفصل \ref{ch:b1:active-margins} --- الهوامش النشطة}

\begin{solution}{exo:b1:active-margins:1}
(أ) $2\pi G\,\Delta\rho\,h = 2\pi\times\num{6.674e-11}\times1700\times3000 =
\qty{2.14e-3}{m.s^{-2}}$: تنخفض الجاذبية بنحو \qty{214}{\milli\gal}. (ب) من أجل \qty{3.5}{km}،
$2\pi\times\num{6.674e-11}\times1700\times3500 = \qty{2.50e-3}{m.s^{-2}}$،
\qty{250}{\milli\gal}: وهو حجم الانخفاض المرصود، فالخندق بعيد جدًا عن توازن أيزوستاتي. (ج) تفترض صيغة الكتلة أن الصخر الغائب يمتد بلا حدود على كل جانب؛ وخندق لا يتجاوز عرضه عشرات الكيلومترات يزيل كتلة أقل بكثير حول نقطة القياس، فماؤه وحده يعطي أقل من \qty{250}{\milli\gal}. ويقع الباقي من العجز أعمق: تنحني الصفيحة كلها نحو الأسفل تحت الخندق، ويحلّ قشرتها، الأخفّ من \omterm{def:b1:layered-earth:layers}{الوشاح}، محلّ \omterm{def:b1:layered-earth:layers}{الوشاح} عند قاعدتها.
\end{solution}

\begin{solution}{exo:b1:active-margins:2}
(أ) $\qty{82}{mm.yr^{-1}} = \qty{82}{km.Myr^{-1}}$، أي \qty{820}{km} في \qty{10}{Myr}. (ب) $820\sin\ang{45} = \qty{580}{km}$. (ج) المسار إلى \qty{600}{km} هو $600/\sin\ang{45} = \qty{849}{km}$، يُقطع في $849/82 = \qty{10.3}{Myr}$.
\end{solution}

\begin{solution}{exo:b1:active-margins:3}
$g = 110/\tan\delta$: \qty{191}{km} عند \ang{30}، و\qty{110}{km} عند \ang{45}، و\qty{64}{km} عند \ang{60}. والقوس فوق الكتلة الأقل انحدارًا يقع أبعد عن خندقه.
\end{solution}

\begin{solution}{exo:b1:active-margins:4}
(أ) التقارب: $\qty{2.7e6}{m^2.yr^{-1}}/\qty{5.1e7}{m} =
\qty{0.053}{m.yr^{-1}}$، أي نحو \qty{53}{mm.yr^{-1}}. والانتشار: $\qty{3.0e6}{m^2.yr^{-1}}/\qty{6.7e7}{m} = \qty{0.045}{m.yr^{-1}}$، أي نحو \qty{45}{mm.yr^{-1}}. على طول كلّي أقصر، تستهلك حدود الاندساس قاع المحيط أسرع مما تصنعه الحيود، لكنها تستهلك أقل في المجموع. (ب) والباقي، نحو \qty{0.3}{km^2}، يُستوعب على طول الحدود المتقاربة الأخرى في النموذج، مثل تصادمات القارات والمناطق المتقاربة داخل المحيطات، حيث تُقصر الصفائح دون أن تندسّ؛ والحيود ليست وحدها حدود التباعد: فالشقوق القارية وبعض التحويلات تنفتح هي أيضًا قليلًا.
\end{solution}

\begin{solution}{exo:b1:active-margins:5}
(أ) $\bar x = 350$، $\bar z = 290$؛ انحرافات $x$ هي $-150$، $-50$،
$50$، $150$، وانحرافات $z$ هي $-140$، $-40$، $40$، $140$، إذن
$S_{xx} = 2(150^2 + 50^2) = \num{50000}$ و
$S_{xz} = 2(150\times140 + 50\times40) = \num{46000}$. ثم $s = 0.92$،
$a = 290 - 0.92\times350 = \qty{-32}{km}$، وانحداره
$\arctan 0.92 = \ang{42.6}$. (ب) $z = 0$ عند $x = 32/0.92 = \qty{34.8}{km}$ من الخندق. (ج) بالنقطتين الطرفيتين يكون الميل $280/300 = 0.933$، أي انحدارًا قدره \ang{43.0}: قريب، لكنه لا يستعمل إلا اثنين من الزلازل الأربعة.
\end{solution}

\begin{solution}{exo:b1:active-margins:6}
(أ) $50/\tan\ang{25} + 60/\tan\ang{45} = 107.2 + 60.0 = \qty{167.2}{km}$. (ب) تعطي كتلة بزاوية \ang{45} على امتدادها كله \qty{110}{km}: فالقطعة الضحلة الأخفّ انحدارًا توسّع الفجوة بمقدار \qty{57}{km}. (ج) القطعة العميقة هي $z = 50 + (x - 107.2)$؛ وتبلغ $z = 0$ عند $x = \qty{57.2}{km}$، نحو اليابسة من الخندق، كما يفعل خط كوريل عند \qty{53}{km}: فقطعة ضحلة أخفّ انحدارًا تفسّر إزاحة كوريل. (د) $107.2 + 60/\tan\ang{60} = 107.2 + 34.6 =
\qty{141.9}{km}$.
\end{solution}

\begin{solution}{exo:b1:active-margins:7}
(أ) نصف القطر الزاوي هو نصف \omterm{def:b1:geological-maps:attitude}{الانحدار}، و$r = R\,\theta$ مع $\theta$ بالراديان: \ang{22.5}، و$6371\times0.3927 = \qty{2502}{km}$ من أجل \ang{45}؛ و\ang{15}، و\qty{1668}{km} من أجل \ang{30}؛ و\ang{35}، و\qty{3892}{km} من أجل \ang{70}. (ب) الأقواس فوق الكتل الأخفّ انحدارًا. (ج) التنبؤ هو $44.9/2 = \ang{22.45}$؛ والتركيب، \ang{21.9}، أصغر منه بمقدار $0.55/22.45 = \qty{2.4}{\percent}$.
\end{solution}

\begin{solution}{exo:b1:active-margins:8}
(أ) $q/k$: $0.035/2.5 = \qty{0.014}{K.m^{-1}} = \qty{14}{K.km^{-1}}$ تحت المنطقة الأمامية، و$0.080/2.5 = \qty{32}{K.km^{-1}}$ تحت \omterm{def:b1:active-margins:active-margin}{حوض خلف القوس}. (ب) $10 + 14\times30 = \qty{430}{\celsius}$ و$10 + 32\times30 = \qty{970}{\celsius}$. (ج) يُنتج جزء من التدفّق الحراري السطحي بالنشاط الإشعاعي داخل \omterm{def:b1:layered-earth:layers}{القشرة} نفسها، فيعبر أقل من الحرارة المستويات الأعمق، ويتناقص التدرّج مع العمق: فالتمديد الخطي لتدرّج السطح يبالغ في تقدير الحرارة عند العمق.
\end{solution}

\begin{solution}{exo:b1:active-margins:9}
(أ) \qty{32}{mm.yr^{-1}} تساوي \qty{32}{km.Myr^{-1}}؛ والمعدّل النصفي \qty{16}{mm.yr^{-1}}. (ب) $100/32 = \qty{3.1}{Myr}$. (ج) لا يلزم أن يكون المعدّل ثابتًا؛ وقد يأتي جزء من العرض من تمدّد القوس وانشقاقه قبل أن يبدأ الانتشار؛ وقد يكون المحور قد قفز أو انتشر بعدم تماثل.
\end{solution}

\begin{solution}{exo:b1:active-margins:10}
(أ) المسار $600/\sin\ang{45} = \qty{849}{km}$، يُقطع بسرعة \qty{82}{km.Myr^{-1}} في \qty{10.3}{Myr}، أي \qty{3.27e14}{s}. (ب) $\sqrt{\num{e-6}\times\num{3.27e14}} = \qty{1.8e4}{m}$، أي نحو \qty{18}{km}. (ج) $(\qty{5e4}{m})^2/\qty{e-6}{m^2.s^{-1}} = \qty{2.5e15}{s}$، أي نحو \qty{79}{Myr}. (د) لم تصل الحرارة إلا إلى \qty{18}{km} الخارجية من كل جانب من الكتلة: فلُبّها يصل إلى \qty{600}{km} باردًا تقريبًا كما غادر الخندق، وأبرد وأصلب من \omterm{def:b1:layered-earth:layers}{الوشاح} المحيط به، ولذلك يمكنه أن ينكسر بالزلازل على ذلك العمق.
\end{solution}

\begin{solution}{exo:b1:active-margins:11}
(أ) على المسافة $L$ خلف طرف الكومة، يقع السطح على ارتفاع $L\tan\alpha$ فوق مستوى الطرف، وتقع القاعدة على عمق $L\tan\beta$ تحته: فخلف الكومة ارتفاعها $L(\tan\alpha + \tan\beta)$، ومساحة المثلث $\frac12 L\times L(\tan\alpha + \tan\beta)$. (ب) $\tan\ang{4} + \tan\ang{6} =
0.175$. وتستقبل الكومة $\qty{1}{km}\times\qty{40}{km.Myr^{-1}} =
\qty{40}{km^2.Myr^{-1}}$ من المقطع، إذن $A = 40t$ و$L = \sqrt{2A/0.175}$: \qty{21.4}{km} بعد \qty{1}{Myr}، و\qty{42.8}{km}
بعد \qty{4}{Myr}، و\qty{67.6}{km} بعد \qty{10}{Myr}. ينمو $L$ بتناسب $\sqrt t$: أربعة أمثال الزمن تعطي ضعف العرض. (ج) تُقسم $A$ على اثنين، فيُقسم $L$ على $\sqrt2$: \qty{15.1}{km}، و\qty{30.2}{km}، و\qty{47.8}{km}.
\end{solution}

\begin{solution}{exo:b1:active-margins:12}
(أ) $\partial F/\partial a = -2\sum(z_i - a - s x_i) = 0$ يعطي
$n\bar z = na + sn\bar x$، أي $\bar z = a + s\bar x$. (ب) يكبر كل $z_i$ و$\bar z$ بمقدار \qty{10}{km}، فالانحرافات $z_i - \bar z$ و$S_{xz}$ لا تتغيّر: فالميل نفسه، و$a = \bar z - s\bar x$ يكبر بمقدار \qty{10}{km}. (ج) انحرافات $x_i - \bar x$ لا تتغيّر، فالميل نفسه؛ و$\bar x$ يقل بمقدار 20، فيزداد $a$ بمقدار $20s$. (د) لا شيء منهما: \omterm{def:b1:geological-maps:attitude}{فالانحدار} لا يتعلّق إلا بالميل. والأخطاء المشتركة بين كل النقاط تُزيح الخط دون أن تُديره.
\end{solution}

\begin{solution}{pb:b1:active-margins:1}
\textbf{1.} تقع النقاط الثماني قرب خط مستقيم ينحدر بنحو \ang{45}، من عمق \qty{96}{km} على بُعد \qty{155}{km} من الخندق إلى عمق \qty{573}{km} عند \qty{609}{km}.

\textbf{2.} $\sum x_i = 3152$, $\sum z_i = 2751$, $\sum x_i^2 = \num{1442910}$,
$\sum x_i z_i = \num{1290548}$; $\bar x = \qty{394.0}{km}$,
$\bar z = \qty{343.875}{km}$.

\textbf{3.} $S_{xx} = 1442910 - 8\times394^2 = \num{201022}$ و$S_{xz} = 1290548 - 8\times394\times343.875 = \num{206654}$. ثم $s = 206654/201022 = 1.028$ و$a = 343.875 - 1.02802\times394 =
\qty{-61.2}{km}$: $z = -61.2 + 1.028\,x$.

\textbf{4.} $\arctan1.028 = \ang{45.8}$.

\textbf{5.} عند $x = 61.2/1.028 = \qty{59.5}{km}$ نحو اليابسة من الخندق. كانت كتلة تغادر الخندق بهذا \omterm{def:b1:geological-maps:attitude}{الانحدار} ستبلغ السطح هناك؛ لكنها تبلغه عند الخندق، فهي تنحدر أخفّ قرب الخندق وتشتدّ تحته.

\textbf{6.} الانحرافات هي $-2.1$، $2.0$، $-7.3$، $5.3$، $7.1$، $2.3$، $-15.0$ و$8.1$ \unit{km}. يقع الزلزال 7 على بعد \qty{15}{km} فوق الخط، والزلزال 8 على بعد \qty{8}{km} تحته؛ وتقع البقية ضمن نحو \qty{7}{km} منه، وهو جزء من سمك الكتلة.

\textbf{7.} $\ang{320.9} - \ang{300} = \ang{20.9}$. عمودية على الخندق: $87.4\cos\ang{20.9} = \qty{81.6}{mm.yr^{-1}}$؛ وموازية له:
$87.4\sin\ang{20.9} = \qty{31.2}{mm.yr^{-1}}$. فالتقارب مائل.

\textbf{8.} \qty{81.6}{mm.yr^{-1}} تساوي \qty{81.6}{km.Myr^{-1}}: \qty{816}{km} في \qty{10}{Myr}.

\textbf{9.} المسار $632/\sin\ang{45} = \qty{894}{km}$، يُقطع في
$894/81.6 = \qty{11.0}{Myr}$.

\textbf{10.} $122 + 11 = \qty{133}{Myr}$.

\textbf{11.} $t = \qty{10.95}{Myr} = \qty{3.46e14}{s}$ و
$\sqrt{\num{e-6}\times\num{3.46e14}} = \qty{1.86e4}{m}$، أي نحو \qty{19}{km} من كل وجه: أقل بكثير من \qty{50}{km} إلى منتصف كتلة سمكها \qty{100}{km}، ولُبّها لا يزال باردًا.

\textbf{12.} $q/k$: \qty{17.2}{K.km^{-1}} تحت قاع المحيط القديم،
\qty{14.0}{K.km^{-1}} تحت المنطقة الأمامية، و\qty{32.0}{K.km^{-1}} تحت \omterm{def:b1:active-margins:active-margin}{حوض خلف القوس}.

\textbf{13.} تحت المنطقة الأمامية تقع الكتلة الباردة على بضعة عشرات الكيلومترات فقط تحت السطح؛ فهي تمتص الحرارة من الصخر فوقها، فتصل حرارة أقل إلى السطح.

\textbf{14.} $h = \qty{2.60e-3}{m.s^{-2}}/(2\pi\times\num{6.674e-11}\times1700)
= \qty{3647}{m}$، أي نحو \qty{3.6}{km}: بقدر، أو أكثر قليلًا، من العمق الكامل للخندق تحت القاع المحيط به، نحو \qty{3.5}{km}. فالخندق ينقصه كتلة لا يعوّضها شيء.

\textbf{15.} فرق الكثافة $3300 - 1000 = \qty{2300}{kg.m^{-3}}$،
و$2\pi\times\num{6.674e-11}\times2300\times3500 = \qty{3.38e-3}{m.s^{-2}}$،
\qty{338}{\milli\gal}. تحسب صيغة الكتلة الكتلة الغائبة إلى أيّ مسافة؛ أما الخندق الحقيقي فينحني نحو الأسفل فقط على شريط عرضه عشرات الكيلومترات، فعجزه أصغر، نحو \qty{260}{\milli\gal}، بين التقديرين الاثنين للكتلة.

\textbf{16.} قوّة الصفيحة المنحنية، التي تُبقي الخندق منخفضًا والرفعة الخارجية مرتفعة كما تفعل صفيحة مرنة، وجذب الكتلة الكثيفة الغائصة في \omterm{def:b1:layered-earth:layers}{الوشاح}.

\textbf{17.} $-61.2 + 1.028\times164 = \qty{107.4}{km}$.

\textbf{18.} نعم: \qty{107}{km} تقع بين 72 و\qty{173}{km}، وهي أعمق من المتوسط، \qty{105}{km}، بنحو \qty{2}{km} فقط.

\textbf{19.} $x = (100 + 61.2)/1.028 = \qty{156.8}{km}$.

\textbf{20.} $\ang{45.8}/2 = \ang{22.9}$، و
$6371\times22.9\times\pi/180 = \qty{2546}{km}$.

\textbf{21.} التركيب، \ang{21.9}، أصغر بمقدار \ang{1.0} تقريبًا، أي نحو \qty{4}{\percent}: \omterm{def:b1:layered-earth:layers}{فالقشرة} المنحنية تصف انحناء قوس كوريل جيدًا.

\textbf{22.} $x = (105 + 61.2)/1.028 = 166.2/1.028 = \qty{161.7}{km}$.

\textbf{23.} $164 - 161.7 = \qty{2.3}{km}$ أقرب إلى الخندق من الجبهة المرصودة، أي $2.3/164 = \qty{1.4}{\percent}$: فبراكين كوريل تقف فوق كتلة أعمق قليلًا جدًا من المتوسط، كما وجد السؤال 18.

\textbf{24.} بأقرب كيلومتر، يتنبأ خط الكتلة الكورية \omterm{def:b1:active-margins:volcanic-arc}{بفجوة القوس والخندق} قدرها \textbf{\qty[text-series-to-math=true]{162}{km}} للبراكين فوق الكتلة على \qty{105}{km} (المرصود: \qty{164}{km}).
\end{solution}
